\Conclusion % заключение к отчёту

По результатам работы решены следующие задачи:
\begin{itemize}
  \item выполнена формальная постановка задачи трёхмерной паллетизации и введены необходимые метрические показатели;
  \item разработаны и программно реализованы четыре эвристики: жадная \emph{row-by-row}, алгоритм «платформ», гибридный прототип и улучшенная послойная укладка;
  \item разработан и внедрён модуль \emph{MAPS}, позволяющий перераспределять коробки между паллетами одного адреса и тем самым повысить среднюю плотность на 23 п.\,п.;
  \item экспериментально определён оптимальный шаг округления габаритов (1 см) и исследовано влияние дискретизации на время работы алгоритмов;
  \item проведено моделирование на 436 реальных наборов: достигнута средняя плотность 85 \% с условием крупной поставке и 67\% иначе при машинном времени 0.019 с на паллету;
\end{itemize}

  Разработанное программное решение реализует четыре алгоритма паллетизации и модуль адресного перераспределения. 
  
Достигнута средняя плотность 62 \%  с СКО $6\%$ при времени вычисления   $\sim0.013$ секунды.

Также для множества паллет в один адрес достигнута средняя плотность 85 \% при времени вычисления $\sim0.012$ секунды на паллету.
 
Разработанные  программы могут быть интегрированы в WMS-систему склада для автоматической генерации схем паллетизации ''на лету''. При суточном потоке 10\,000 паллет время расчёта не превышает трёх минут, что сопоставимо с временем обмена данными между подсистемами.
 

В дальнейшем планируется исследовать обучение эвристики при помощи методов RL для on-line-сценариев и внедрить данный алгоритм в существующие WMS. 

